\documentclass[11pt,a4paper]{article}
\usepackage[utf8]{inputenc}
\usepackage[T1]{fontenc}
\usepackage[french]{babel}
\usepackage{lmodern,geometry,microtype}
\usepackage{amsmath}
\usepackage{siunitx}
\usepackage[version=4]{mhchem}
\usepackage{xcolor,listings,hyperref}
\geometry{margin=2.4cm}
\sisetup{locale=FR,per-mode=symbol}
\DeclareSIUnit\ppm{ppm}
\lstset{language=Python,basicstyle=\ttfamily\small,backgroundcolor=\color{black!4},frame=single,breaklines=true,showstringspaces=false}
\hypersetup{colorlinks=true,linkcolor=blue!45!black,urlcolor=blue!55!black}

\title{Documentation de \texttt{code\_H20\_CH4\_CO2\_O3.py}}
\author{Projet climat 2026}
\date{\today}

\begin{document}
\maketitle
\tableofcontents

\section{Emplacement et rôle}

Le fichier documenté se trouve dans
\path{../../temperatures_et_concentrations/code_H20_CH4_CO2_O3.py}, relativement
au dossier du modèle HITRAN. Il fournit le profil vertical simplifié utilisé
par \path{code_final_emissivite_avec_sections_hitran.py} : température,
pression et concentrations de \ce{CO2}, \ce{CH4}, \ce{O3} et \ce{H2O}.

Malgré la chaîne \texttt{H20} présente dans son nom, le gaz traité par le
code et par les clés de sortie est bien \texttt{H2O}.

\section{Import préalable}

\begin{lstlisting}
import math
\end{lstlisting}

\texttt{math} appartient à la bibliothèque standard de Python. Aucun paquet
\texttt{pip} supplémentaire n'est requis par ce fichier. Le module est chargé
par \texttt{exec} dans le modèle actuel, qui récupère ensuite la fonction
\texttt{get\_gases\_ppm}.

\section{Fonction \texttt{temperature\_atmosphere(z)}}

\textbf{Entrée :} altitude \texttt{z} en kilomètres.

\textbf{Retour :} température en degrés Celsius.

Le profil est affine par morceaux, avec des limites à 0, 11, 20, 32, 50,
71 et \SI{85}{\kilo\metre}. Les gradients successifs sont $-6{,}5$, 0, 1,
$2{,}8$, $-2{,}8$, $-2$ et \SI{12}{\kelvin\per\kilo\metre}. La dernière
expression est aussi utilisée hors de la plage principale.

\section{Fonction \texttt{get\_gases\_ppm(z\_km)}}

\textbf{Entrée :} altitude en kilomètres.

\textbf{Retour :} dictionnaire contenant altitude, température, pression et
deux grandeurs par gaz : proportion de mélange et quantité volumique absolue.
Le modèle HITRAN utilise les champs de mélange comme concentrations en ppm.

\subsection{Pression atmosphérique}

Les constantes sont $g=\SI{9.80665}{\metre\per\second\squared}$,
$M=\SI{0.0289644}{\kilogram\per\mole}$,
$R=\SI{8.31445}{\joule\per\mole\per\kelvin}$ et
$p_0=\SI{101325}{\pascal}$. Pour une couche à gradient thermique non nul,
\begin{equation}
  p=p_0\left(\frac{T}{T_0}\right)^{-gM/(R\alpha)}.
\end{equation}
Pour une couche isotherme,
\begin{equation}
  p=p_0\exp\left(-\frac{gM\Delta z}{RT_0}\right).
\end{equation}
Les pressions aux limites sont d'abord propagées, puis la formule de la couche
contenant l'altitude demandée est appliquée.

\subsection{Profils des gaz}

\begin{itemize}
  \item \ce{CO2} : proportion constante de \SI{425}{\ppm} ;
  \item \ce{CH4} : \SI{1.9}{\ppm} jusqu'à \SI{11}{\kilo\metre}, puis
        décroissance exponentielle de hauteur \SI{25}{\kilo\metre} ;
  \item \ce{O3} : fond de \SI{0.1}{\ppm}, avec un profil gaussien entre 15 et
        \SI{35}{\kilo\metre}, centré à \SI{25}{\kilo\metre} et culminant à
        \SI{8}{\ppm} ;
  \item \ce{H2O} : pression de vapeur saturante et humidité relative
        décroissante sous \SI{11}{\kilo\metre}, avec un plancher de
        \SI{4}{\ppm}.
\end{itemize}

Les quantités dites « volumique absolu » sont corrigées par
$(p/p_{\mathrm{sol}})(T_{\mathrm{sol}}/T)$, mais elles ne sont pas utilisées
par le calcul HITRAN actuel.

\section{Clés du dictionnaire utiles au modèle HITRAN}

\begin{lstlisting}
"Temperature (°C)"
"Pression locale (Pa)"
"CO2 de mélange (Proportion)"
"H2O de mélange (Proportion)"
"CH4 de mélange (Proportion)"
"O3 de mélange (Proportion)"
\end{lstlisting}

Toute modification de ces libellés doit être reportée dans
\path{code_final_emissivite_avec_sections_hitran.py} et
\path{generer_table_emissivite_hitran.py}.

\section{Limites}

Le profil est une atmosphère de référence paramétrique, non une observation
météorologique. Les valeurs sont arrondies dans le dictionnaire. L'extrapolation
thermique au-dessus de \SI{85}{\kilo\metre}, le plancher de vapeur d'eau et le
profil d'ozone sont des simplifications à prendre en compte dans l'analyse des
résultats.

\end{document}
